\chapter{The Klein four-group of characters}
\label{chap:klein}

\section{Preliminaries}
\supporting

The patterns of the characters at a prime along~$(c_n)$. For a prime~$p$ and an integer
$m \ne 0$, $m_p$ is the prime-to-$p$ part of~$m$, and ``$m \ne \pm\sq$'' means that~$|m|$ is not
a square.

\begin{lemma}[Square classes of the terms]
  \label{lem:sqclass}
  \uses{lem:reduction, def:qchar}
  \lean{QuadraticIterates.mulChar_cSeq_one, QuadraticIterates.mulChar_cSeq_of_two_le,
    QuadraticIterates.mulChar_cSeq_one_eq_mul}
  \leanok
  For every quadratic character~$\chi$: $\chi(c_1) = \chi(-rs)$ and $\chi(c_n) = \chi(r w_n)$ for
  $n \ge 2$. In particular $\chi(c_2) = \chi(r(r+s))$.
\end{lemma}
\begin{proof}
  \leanok
  $c_n = r w_n / s^{2^{n-1}}$ (Lemma~\ref{lem:reduction}), whose denominator is a square for
  $n \ge 2$, and $c_1 = -rs/s^2$; a character takes the same value at~$x$ and at~$x y^2$.
\end{proof}

\begin{lemma}[Two-periodicity modulo $(r + s)^2$]
  \label{lem:strongcong}
  \uses{lem:reduction, lem:wseq-mul}
  \lean{QuadraticIterates.wSeq_two_mul_modEq, QuadraticIterates.wSeq_two_mul_add_one_modEq,
    QuadraticIterates.exists_wSeq_two_mul_eq, QuadraticIterates.mul_wSeq_two_mul_add_one_modEq,
    QuadraticIterates.cSeq_two_mul_of_sq_eq_zero,
    QuadraticIterates.cSeq_two_mul_add_one_of_sq_eq_zero}
  \leanok
  For $k \ge 1$: $w_{2k} = (r + s)\, T_k$ with $T_k \equiv s^{2^{2k-1} - 2} \bmod (r + s)$, and
  $w_{2k+1} \equiv s^{2^{2k} - 1} \bmod (r + s)$.
  \formalnote{The first part is the case $m = 2$ of Lemma~\ref{lem:wseq-mul}, as $w_2 = r + s$,
  stated as $w_{2k} \equiv (r + s) s^{2^{2k-1} - 2} \bmod (r + s)^2$. The second uses the
  general fact that $\bar c_{2k} = -t$ and $\bar c_{2k+1} = t - 1$ for $k \ge 1$ in a ring with
  $t^2 = 0$, $t = \bar a + 1$.}
\end{lemma}
\begin{proof}
  \leanok
  Modulo~$r + s$, $s$ is a unit and $\bar a = -1$, so $(\bar c_n) = (1, 0, -1, 0, -1, \dots)$, and
  Lemma~\ref{lem:reduction} gives $r w_{2k+1} \equiv -s^{2^{2k}}$, where $r \equiv -s$. The first
  part is Lemma~\ref{lem:wseq-mul} with $m = 2$.
\end{proof}

\begin{lemma}[Primes of~$r$]
  \label{lem:pat-r}
  \uses{lem:sqclass, lem:wseq-res, def:chars, lem:vp}
  \lean{QuadraticIterates.valChar_cSeq_of_dvd_left,
    QuadraticIterates.unitPartChar_cSeq_of_dvd_left}
  \leanok
  Let~$p$ be a prime dividing~$r$.
  \begin{enumerate}
    \item $v_p(c_n) = v_p(r)$ for all $n \ge 1$; if $v_p(r)$ is odd, $\val_p$ is a~$\phi_0$.
    \item If~$p$ is odd and $x = \Leg{r_p s}{p}$, then
          $\chi_p(c_\bullet) = (\Leg{-1}{p} x, x, x, \dots)$.
    \item If $4 \mid r$ and $x = \chi_4(r_2 s)$, then $\chi_4(c_\bullet) = (-x, x, x, \dots)$.
    \item If $8 \mid r$, then $\chi_8(c_n) = \chi_8(r_2 s)$ for all $n \ge 1$.
  \end{enumerate}
  \formalnote{(1) is \lname{padicValRat\_cSeq\_of\_dvd\_left} (Lemma~\ref{lem:vp}), and its
  parity \lname{valChar\_cSeq\_of\_dvd\_left}. (2)--(4) are one statement,
  \lname{unitPartChar\_cSeq\_of\_dvd\_left}: for a power $N = p^k \mid r$, $k \ge 1$, and~$\psi$
  modulo~$N$, $\chi_\psi(c_n) = \psi(s)\chi_\psi(r)$ for $n \ge 2$; at $n = 1$,
  $\chi(c_1) = \chi(-1)\chi(r)\chi(s)$ for every character~$\chi$
  (\lname{mulChar\_cSeq\_one\_eq\_mul}).}
\end{lemma}
\begin{proof}
  \leanok
  By Lemma~\ref{lem:sqclass}, $\chi(c_n) = \chi(r)\chi(w_n)$ for $n \ge 2$ and
  $\chi(c_1) = \chi(-1)\chi(r)\chi(s)$. As $w_n \equiv s^{2^{n-1}-1} \bmod r$ with an odd
  exponent, $\chi_\psi(w_n) = \psi(s)$ for~$\psi$ modulo a divisor~$N$ of~$r$; and
  $\chi_\psi(r) = \psi(r_p)$. For (1), $v_p(w_n) = v_p(s) = 0$.
\end{proof}

\begin{lemma}[Primes of~$s$]
  \label{lem:pat-s}
  \uses{lem:sqclass, lem:wseq-res, def:chars}
  \lean{QuadraticIterates.valChar_cSeq_of_dvd_right,
    QuadraticIterates.unitPartChar_cSeq_of_dvd_right,
    QuadraticIterates.unitPartChar_two_cSeq_of_even}
  \leanok
  \begin{enumerate}
    \item Let~$p$ be an odd prime dividing~$s$. Then $\val_p(c_n) = \chi_p(c_n) = 1$ for $n \ge 2$,
          while $\val_p(c_1) = (-1)^{v_p(s)}$ and $\chi_p(c_1) = \Leg{-r s_p}{p}$.
    \item Let~$s$ be even. Then $\val_2(c_n) = \chi_4(c_n) = \chi_8(c_n) = 1$ for $n \ge 3$, and
          $\val_2(c_2) = 1$.
  \end{enumerate}
  In particular, $c_n$ is a square in~$\QQ_p$ for $n \ge 2$, resp.\ in~$\QQ_2$ for $n \ge 3$.
  \formalnote{(1) for~$\psi$ modulo any power $N \mid s$ of~$p$, $k \ge 1$; (2) for~$\psi$
  modulo a power of~$2$ dividing~$8$. The values at~$c_1$ are those of
  \lname{mulChar\_cSeq\_one\_eq\_mul}.}
\end{lemma}
\begin{proof}
  \leanok
  For $n \ge 2$, $v_p(c_n) = -2^{n-1} v_p(s)$ is even. (1) $w_n \equiv r^{2^{n-1}-1} \bmod p$, so
  $r w_n \equiv r^{2^{n-1}}$ is a nonzero square modulo~$p$. (2) $w_n \equiv r \bmod 8$ for
  $n \ge 3$ (Lemma~\ref{lem:wseq-res}), so $r w_n \equiv r^2 \equiv 1 \bmod 8$.
\end{proof}

\begin{lemma}[Primes of~$r + s$]
  \label{lem:pat-rs}
  \uses{lem:sqclass, lem:strongcong, def:chars}
  \lean{QuadraticIterates.valChar_cSeq_two_mul_add_one_of_dvd_add,
    QuadraticIterates.valChar_cSeq_two_mul_of_dvd_add,
    QuadraticIterates.unitPartChar_cSeq_one_of_dvd_add,
    QuadraticIterates.unitPartChar_cSeq_two_mul_add_one_of_dvd_add,
    QuadraticIterates.unitPartChar_cSeq_two_mul_of_dvd_add}
  \leanok
  \begin{enumerate}
    \item Let~$q$ be a prime dividing~$r + s$. Then $v_q(c_n) = 0$ for odd~$n$ and
          $v_q(c_n) = v_q(r + s)$ for even~$n$; so if $v_q(r + s)$ is odd,
          $\val_q(c_\bullet) = ({+}{-}{+}{-}\cdots)$.
    \item Let $p \equiv 3 \bmod 4$ be a prime dividing~$r + s$, and $\chi = \chi_p$, or let
          $4 \mid r + s$, $p = 2$ and $\chi = \chi_4$. Then
          \[ \chi(c_\bullet) = (1, y, -1, y, -1, y, \dots), \qquad
             y = \chi(c_2) = -\chi(s\,(r + s)_p). \]
  \end{enumerate}
  \formalnote{(2) is stated for~$\psi$ modulo any power $N = p^k \mid r + s$, $k \ge 1$:
  $\chi_\psi(c_1) = 1$, $\chi_\psi(c_{2j+1}) = \psi(-1)$ and
  $\chi_\psi(c_{2j}) = \psi(-1)\psi(s)\chi_\psi(r + s)$ for $j \ge 1$.}
\end{lemma}
\begin{proof}
  \leanok
  By Lemma~\ref{lem:strongcong}, $w_{2k} = (r + s) T_k$ with $T_k \equiv s^{2^{2k-1} - 2}$ and
  $w_{2k+1} \equiv s^{2^{2k} - 1} \bmod (r + s)$ for $k \ge 1$, and $q \nmid rs$. This gives (1),
  as $v_q(c_n) = v_q(w_n)$ and~$s$ is prime to~$q$. For (2): $\chi(c_1) = \chi(-rs) = 1$ because
  $-rs \equiv s^2 \bmod (r + s)$. For odd $n \ge 3$, $\chi(c_n) = \chi(r w_n)$ and
  $r w_n \equiv -s \cdot s^{2^{n-1} - 1} = -s^{2^{n-1}} \bmod (r + s)$, so $\chi(c_n) = \psi(-1) =
  -1$. For even $n = 2k$, $\chi(c_n) = \chi(r (r + s))\psi(T_k) = \chi(r(r+s)) = \chi(c_2)$, the
  exponent $2^{2k-1} - 2$ being even. Finally $\chi(r) = \psi(-s)$ and
  $\chi(r + s) = \psi((r + s)_p)$.
\end{proof}

\begin{lemma}[Primes of~$r + 2s$]
  \label{lem:pat-r2s}
  \uses{lem:localeval, lem:rigid}
  \lean{QuadraticIterates.unitPartChar_cSeq_of_dvd_add_two_mul, QuadraticIterates.cSeq_neg_two}
  \leanok
  Let~$p$ be an odd prime dividing~$r + 2s$. Then $\chi_p(c_n) = \Leg{2}{p}$ for all $n \ge 1$.
  \formalnote{For~$\psi$ modulo any power $N = p^k \mid r + 2s$, $k \ge 1$:
  $\chi_\psi(c_n) = \psi(2)$.}
\end{lemma}
\begin{proof}
  \leanok
  $p \nmid s$, and $\bar a = -2$ in~$\ZZ/p\ZZ$, so $\bar c_1 = 2$ is a fixed point of
  $x \mapsto x^2 - 2$ (so the sequence is constant, Lemma~\ref{lem:rigid}); apply
  Lemma~\ref{lem:localeval}.
\end{proof}

\begin{lemma}[Primes of $r^2 + s^2$]
  \label{lem:pat-r2s2}
  \uses{lem:localeval, lem:rigid}
  \lean{QuadraticIterates.unitPartChar_cSeq_two_mul_add_one_of_dvd_sq_add_sq,
    QuadraticIterates.unitPartChar_cSeq_two_mul_of_dvd_sq_add_sq,
    QuadraticIterates.cSeq_two_mul_add_one_of_sq_eq_neg_one,
    QuadraticIterates.cSeq_two_mul_of_sq_eq_neg_one}
  \leanok
  Let~$p$ be an odd prime dividing $r^2 + s^2$. Then $\chi_p(c_n) = \Leg{-rs}{p}$ for odd~$n$ and
  $\chi_p(c_n) = \Leg{(r - s)s}{p}$ for even~$n$.
  \formalnote{For~$\psi$ modulo any power $N = p^k \mid r^2 + s^2$, $k \ge 1$, with values
  $\psi(-rs)$ and $\psi((r - s)s)$.}
\end{lemma}
\begin{proof}
  \leanok
  $p \nmid rs$ and $\bar a^2 = -1$ in~$\ZZ/p\ZZ$. Since
  $\bar c_3 - \bar c_1 = \bar a(\bar a + 2)(\bar a^2 + 1) = 0$, the sequence has period~$2$ from
  $n = 1$ on, periodicity propagating along the recursion (Lemma~\ref{lem:rigid});
  $\bar c_1 = -\bar a$ and $\bar c_2 = \bar a^2 + \bar a =
  \bar a - 1$ are units ($\bar a = 1$ would give $2 = 0$). Lemma~\ref{lem:localeval} gives
  $\chi_p(c_1) = \Leg{-r s^{-1}}{p} = \Leg{-rs}{p}$ and $\chi_p(c_2) = \Leg{(r-s)s}{p}$.
\end{proof}

\begin{remark}
  As~$\bar a$ has order~$4$ in $(\ZZ/p\ZZ)^\times$, $p \equiv 1 \bmod 4$, and
  $\Leg{-rs}{p} = \Leg{\bar a}{p} = -1$ if and only if $p \equiv 5 \bmod 8$. The direction used
  for (Z3) is \lname{ZMod.quadraticChar\_eq\_neg\_one\_of\_sq\_eq\_neg\_one}.
\end{remark}

We write $a \equiv m \bmod N$ when~$N$ divides the numerator of~$a - m$; for $a \equiv 1 \bmod 4$
this means that~$s$ is odd and $r \equiv s \bmod 4$.

\begin{lemma}[$a \equiv 1 \bmod 4$]
  \label{lem:pat-dyadic}
  \uses{lem:localeval, lem:pat-rs, def:phin}
  \lean{QuadraticIterates.hasCharPattern_mulSingle_one_of_modEq_four,
    QuadraticIterates.valChar_two_cSeq_succ_of_modEq_four,
    QuadraticIterates.unitPartChar_two_χ₄_cSeq_one_of_modEq_four,
    QuadraticIterates.unitPartChar_two_χ₄_cSeq_two_mul_add_one_of_modEq_four,
    QuadraticIterates.unitPartChar_two_χ₄_cSeq_two_mul_of_modEq_four,
    QuadraticIterates.intCast_mul_wSeq_zmod_eight, QuadraticIterates.cSeq_two_mul_zmod_eight,
    QuadraticIterates.cSeq_two_mul_add_one_zmod_eight}
  \leanok
  Let $a \equiv 1 \bmod 4$. Then $\val_2(c_\bullet) = ({+}{-}{+}{-}\cdots)$, and~$\chi_4$ (if
  $a \equiv 1 \bmod 8$) or~$\chi_4 \val_2$ (if $a \equiv 5 \bmod 8$) has pattern
  $({-}{+}{+}\cdots)$: $a$ has~$\phi_1$.
  \formalnote{$a \equiv 1 \bmod 4$ is stated as~$s$ odd and $r \equiv s \bmod 4$. Modulo~$8$,
  $\bar a = rs \in \{1, 5\}$, the even terms are~$\bar a + 1$ and the odd terms from $n = 3$ on
  are~$\bar a + 4$ (computations by \lname{decide}); in all, $\chi_4(c_\bullet) = (-1, w, 1, w, 1,
  \dots)$ with $w = \chi_8(rs)$.}
\end{lemma}
\begin{proof}
  \leanok
  \uses{lem:reduction, lem:orbit}
  $r + s \equiv 2 \bmod 4$, so the statement on~$\val_2$ is Lemma~\ref{lem:pat-rs}(1) with $q = 2$.
  In~$\ZZ/8\ZZ$, $\bar a \in \{1, 5\}$ and $(\bar c_n)$ is $(-1, 2, 5, 2, 5, \dots)$ for
  $\bar a = 1$ and $(3, 6, 1, 6, 1, \dots)$ for $\bar a = 5$. For $n \ge 2$,
  $s^{2^{n-1}} \equiv 1 \bmod 8$, so $r w_n \equiv \bar c_n \bmod 8$ (Lemma~\ref{lem:reduction});
  for even~$n$ the odd part $r w_n / 2$ is $\equiv 1$ resp.\ $3 \bmod 4$. And
  $c_1 \sim -rs \equiv -\bar a \bmod 8$. So $\chi_4(c_\bullet) = ({-}{+}{+}{+}\cdots)$ for
  $\bar a = 1$ and $({-}{-}{+}{-}{+}\cdots)$ for $\bar a = 5$, and~$\val_2$ turns the latter into
  $({-}{+}{+}\cdots)$.
\end{proof}

\section{Criteria for \texorpdfstring{$\phi_0$, $\phi_1$ and~$\tau$}{phi\_0, phi\_1 and tau}}
\paperitems

Lemma~\ref{lem:twoadic} tells us that for half of the rational $a \in \cA$, $\rho_a$ is
surjective \emph{if} we can also produce~$\phi_0$ and~$\phi_1$. This is equivalent to~$a$ having
two of~$\phi_0$, $\phi_1$ and~$\tau$: if~$\phi_0$ and~$\phi_1$ are such characters, then
$\{1, \phi_0, \phi_1, \phi_0\phi_1\}$ is a Klein four-group inside the character group of
$\QQ^\times/\QQ^{\times 2}$, and~$\phi_0\phi_1$ is a~$\tau$ (Lemma~\ref{lem:klein}). We now collect
sufficient conditions for each of them. For a quadratic character~$\psi$ on~$\QQ^\times$, we have
$\psi(c_1) = \psi(-rs)$, $\psi(c_2) = \psi(r(r + s))$ and $\psi(c_1 c_2) = \psi(-s(r + s))$
(Lemma~\ref{lem:sqclass}).

\begin{proposition}[Criteria for~$\phi_0$]
  \label{prop:Z}
  \uses{def:phi0tau, lem:pat-r, lem:pat-r2s, lem:pat-r2s2}
  \lean{QuadraticIterates.NegAllCriterion,
    QuadraticIterates.hasCharPattern_neg_one_of_negAllCriterion,
    QuadraticIterates.hasCharPattern_neg_one_of_not_isSquare,
    QuadraticIterates.hasCharPattern_neg_one_of_dvd_add_two_mul,
    QuadraticIterates.hasCharPattern_neg_one_of_dvd_sq_add_sq,
    ZMod.quadraticChar_eq_neg_one_of_sq_eq_neg_one,
    QuadraticIterates.hasCharPattern_neg_one_of_dvd_left,
    QuadraticIterates.hasCharPattern_neg_one_of_dvd_left_two}
  \leanok
  $a$ has~$\phi_0$ if one of the following holds.
  \begin{enumerate}
    \item[(Z1)] $|r|$ is not a square.
    \item[(Z2)] A prime $p \equiv \pm 3 \bmod 8$ divides~$r + 2s$.
    \item[(Z3)] A prime $p \equiv 5 \bmod 8$ divides $r^2 + s^2$ with
          $\chi_p\bigl((r - s)s\bigr) = -1$.
    \item[(Z4)] A prime $p \equiv 1 \bmod 4$ divides~$r$ with $\chi_p(-rs) = -1$; or $8 \mid r$
          and $\chi_8(-rs) = -1$.
  \end{enumerate}
  \formalnote{An inductive proposition \lname{NegAllCriterion r s} with one constructor per item
  (two for (Z4)), stated with Mathlib's notions only: Legendre symbols as \lname{jacobiSym}
  and~$p$-parts as $r = p^k m$ with $p \nmid m$. In (Z4), $\chi_p(-rs) = -1$ is
  $\Leg{r_p s}{p} = -1$, as $\chi_p(-1) = 1$, and $\chi_8(-rs) = -1$ is
  $r_2 s \equiv \pm 3 \bmod 8$. In (Z3), $p \nmid s$ follows from
  $\Leg{(r - s)s}{p} = -1$, and $\Leg{-rs}{p} = -1$ is
  \lname{ZMod.quadraticChar\_eq\_neg\_one\_of\_sq\_eq\_neg\_one} (Euler's criterion).}
\end{proposition}
\begin{proof}
  \leanok
  (Z1) $v_p(c_n) = v_p(r)$ for all~$n$, so~$\val_p$ for a prime with $v_p(r)$ odd is a suitable
  character (Lemma~\ref{lem:pat-r}(1)). (Z2) $c_n \equiv 2 \bmod p$ for all~$n$, so we can
  take~$\chi_p$ (Lemma~\ref{lem:pat-r2s}). (Z3) Here $a^2 \equiv -1 \bmod p$, so the sequence $(c_n)_{n \ge 1}$ reduces
  modulo~$p$ to $(-a, a - 1, -a, a - 1, \dots)$ (Lemma~\ref{lem:pat-r2s2}); moreover
  $\chi_p(-a) = \chi_p(a) = (-1)^{(p-1)/4} = -1$ and $\chi_p(a - 1) = \chi_p((r - s)s) = -1$. So
  we can take~$\chi_p$. (Z4)
  $\chi_p(c_1) = \chi_p(-1)\chi_p(rs)$ and $\chi_p(c_n) = \chi_p(rs)$ for $n \ge 2$, and
  similarly for~$\chi_8$ in the second case (Lemma~\ref{lem:pat-r}(2), (4)); so we can
  take~$\chi_p$, resp.\ $\chi_8$.
\end{proof}

\begin{proposition}[Criteria for~$\phi_1$]
  \label{prop:O}
  \uses{lem:signs, lem:pat-dyadic, lem:pat-s, lem:pat-r}
  \lean{QuadraticIterates.NegFirstCriterion,
    QuadraticIterates.hasCharPattern_mulSingle_one_of_negFirstCriterion,
    QuadraticIterates.hasCharPattern_mulSingle_one_of_dvd_right_of_odd,
    QuadraticIterates.hasCharPattern_mulSingle_one_of_dvd_right_of_apply,
    QuadraticIterates.hasCharPattern_mulSingle_one_of_dvd_right,
    QuadraticIterates.hasCharPattern_mulSingle_one_of_dvd_right_two,
    QuadraticIterates.hasCharPattern_mulSingle_one_of_dvd_right_four,
    QuadraticIterates.hasCharPattern_mulSingle_one_of_dvd_right_eight,
    QuadraticIterates.hasCharPattern_mulSingle_one_of_dvd_left,
    QuadraticIterates.hasCharPattern_mulSingle_one_of_dvd_left_four}
  \leanok
  $a$ has~$\phi_1$ if one of the following holds.
  \begin{enumerate}
    \item[(O1)] $a > 0$.
    \item[(O2)] $a \equiv 1 \bmod 4$, i.e.\ $s$ is odd and $r \equiv s \bmod 4$.
    \item[(O3)] For some prime~$p$ dividing~$s$, $-rs$ is not a square in~$\QQ_p$, and if $p = 2$
          and $v_2(s) = 2$, then moreover $\chi_4(-rs) = -1$.
    \item[(O4)] A prime $p \equiv 3 \bmod 4$ divides~$r$ with $\chi_p(-rs) = -1$; or $4 \mid r$
          and $\chi_4(-rs) = -1$.
  \end{enumerate}
  In particular, by (O3), $a$ has~$\phi_1$ if~$s$ is not a square.
  \formalnote{\lname{NegFirstCriterion r s}. (O3) has one constructor at an odd prime, with
  $v_p(s)$ odd or $\Leg{-r s_p}{p} = -1$, and one at~$2$, with $v_2(s)$ odd, or $v_2(s) = 2$ and
  $r s_2 \equiv 1 \bmod 4$, or $v_2(s) \ge 4$ even and $r s_2 \not\equiv -1 \bmod 8$: these are
  the explicit forms of ``$-rs$ is not a square in~$\QQ_p$''. At~$2$, $\chi_4$ works whenever
  $4 \mid s$ and $r s_2 \equiv 1 \bmod 4$, and~$\chi_8$ whenever $8 \mid s$ and
  $r s_2 \equiv \pm 3 \bmod 8$. In (O4), $\chi_p(-rs) = -1$ is $\Leg{r_p s}{p} = 1$, as
  $\chi_p(-1) = -1$, and $\chi_4(-rs) = -1$ is $r_2 s \equiv 1 \bmod 4$.}
\end{proposition}
\begin{proof}
  \leanok
  \uses{lem:sqclass}
  (O1) The sign (Lemma~\ref{lem:signs}). (O2) $\chi_4$ or~$\chi_4 \val_2$, by a consideration
  modulo~$8$; note $\val_2(c_n) = (-1)^{n-1}$ (Lemma~\ref{lem:pat-dyadic}). (O3) Note that
  $r w_n \equiv r^{2^{n-1}} \bmod s$ for $n \ge 2$ (Lemma~\ref{lem:pat-s}(1)). If $v_p(s)$ is odd,
  we can take~$\val_p$. If $v_p(s) \ge 2$ is even and~$p$ is odd, we can take~$\chi_p$: the
  assumption on~$-rs$ gives $\chi_p(c_1) = -1$, and all~$c_n$ with $n \ge 2$ are~$p$-adic squares.
  If $v_2(s) \ge 2$ is even, then by Lemma~\ref{lem:pat-s}(2) the characters at~$2$ are trivial
  on~$c_n$ for $n \ge 3$; here $c_1 \sim -rs$, $c_2 \sim r(r+s)$ (Lemma~\ref{lem:sqclass}), $r$
  odd. If $v_2(s) = 2$, then $r + s \equiv r + 4 \bmod 8$, so $\chi_4(c_2) = \chi_4(r^2) = 1$, and
  $\chi_4(c_1) = \chi_4(-rs) = -1$ by the additional assumption. If $v_2(s) \ge 4$, then
  $r(r + s) \equiv 1 \bmod 8$, so every character at~$2$ is~$1$ at~$c_2$, and as~$-rs$ is not a
  square in~$\QQ_2$, one of~$\chi_4$, $\chi_8$, $\chi_4\chi_8$ is~$-1$ at~$c_1$. (O4) $\chi_p$,
  resp.\ $\chi_4$ (Lemma~\ref{lem:pat-r}(2), (3)); compare the proof of (Z4).
\end{proof}

\begin{proposition}[Criteria for~$\tau$]
  \label{prop:T}
  \uses{lem:signs, lem:window, lem:pat-r, lem:pat-rs}
  \lean{QuadraticIterates.NegTailCriterion,
    QuadraticIterates.hasCharPattern_neg_mulSingle_one_of_negTailCriterion,
    QuadraticIterates.hasCharPattern_neg_mulSingle_one_of_neg_lt,
    QuadraticIterates.hasCharPattern_neg_mulSingle_one_of_dvd_left,
    QuadraticIterates.hasCharPattern_neg_mulSingle_one_of_dvd_left_four,
    QuadraticIterates.hasCharPattern_neg_mulSingle_one_of_dvd_add,
    QuadraticIterates.hasCharPattern_neg_mulSingle_one_of_dvd_add_of_valChar,
    QuadraticIterates.hasCharPattern_neg_mulSingle_one_of_dvd_add_prime,
    QuadraticIterates.hasCharPattern_neg_mulSingle_one_of_dvd_add_four,
    QuadraticIterates.hasCharPattern_neg_mulSingle_one_of_not_isSquare_add,
    QuadraticIterates.hasCharPattern_neg_mulSingle_one_of_window,
    QuadraticIterates.mul_wSeq_three, QuadraticIterates.mul_wSeq_four}
  \leanok
  $a$ has~$\tau$ if one of the following holds.
  \begin{enumerate}
    \item[(T1)] $-1 < a < 0$.
    \item[(T2)] A prime $p \equiv 3 \bmod 4$ divides~$r$ with $\chi_p(-rs) = 1$; or $4 \mid r$
          and $\chi_4(-rs) = 1$.
    \item[(T3)] A prime $p \equiv 3 \bmod 4$ divides~$r + s$ with $\chi_p(-s(r + s)) = -1$; or
          $4 \mid r + s$ and $\chi_4(-s(r + s)) = -1$; or $|r + s|$ is not a square and
          divisible by~$4$ or by a prime $p \equiv 3 \bmod 4$.
    \item[(T4)] $\alpha' < a < -1$ (that is, $c_3 < 0 < c_4$), and $|r + s|$ is not a square.
  \end{enumerate}
  \formalnote{\lname{NegTailCriterion r s}. As $\chi_p(-1) = -1$, the conditions of (T2) and (T3)
  are $\Leg{r_p s}{p} = -1$, $r_2 s \equiv 3 \bmod 4$, $\Leg{s (r+s)_p}{p} = 1$ and
  $s (r+s)_2 \equiv 1 \bmod 4$. In (T4) the window is stated as $r < -s$ and
  $r w_3 < 0 < r w_4$, with $r w_3 = (r^2 + rs)^2 + rs^3$ and $r w_4 = (r w_3)^2 + rs^7$.}
\end{proposition}
\begin{proof}
  \leanok
  \uses{lem:pat-rs}
  (T1) The sign (Lemma~\ref{lem:signs}). (T2) $\chi_p$ resp.\ $\chi_4$
  (Lemma~\ref{lem:pat-r}(2), (3)). (T3) Let $\chi = \chi_p$, resp.\ $\chi_4$. Here $a \equiv -1
  \bmod p$, so modulo~$p$, $(c_n)_{n \ge 1}$ reduces to $(1, 0, -1, 0, -1, \dots)$, and
  $(c_{2n})_p \equiv (c_2)_p \bmod p$, as $c_{2k} \equiv c_2 \bmod c_2^2$ (Lemma~\ref{lem:pat-rs}).
  So~$\chi$ is constant on $(c_{2n})$; take~$\chi$ when $\chi(c_2) = -1$ and~$\chi \val_q$ for a
  prime~$q$ with $v_q(r + s)$ odd otherwise. (T4) $\sgn \cdot \val_q$ with $v_q(r + s)$ odd:
  $({+}{+}{-}{+}{-}\cdots) \cdot ({+}{-}{+}{-}\cdots) = ({+}{-}{-}{-}\cdots)$
  (Lemmas~\ref{lem:window} and~\ref{lem:pat-rs}(1)).
\end{proof}

We say that ``$a$ satisfies (Z)'', ``$a$ satisfies (O)'' and ``$a$ satisfies (T)'' when~$a$
satisfies one of the conditions (Z1)--(Z4), one of (O1)--(O4), and one of (T1)--(T4),
respectively.

We note a useful consequence of (Z1), (Z2) and (O2).

\begin{lemma}[$a \equiv 1 \bmod 4$]
  \label{lem:mod4}
  \uses{prop:Z, prop:O, lem:prime3}
  \lean{QuadraticIterates.negAllCriterion_of_modEq_four,
    QuadraticIterates.hasCharPattern_neg_one_of_modEq_four}
  \leanok
  If $a \equiv 1 \bmod 4$, then~$a$ has~$\phi_0$ and~$\phi_1$.
  \formalnote{The condition $a \equiv 1 \bmod 4$ is stated as~$s$ odd and $r \equiv s \bmod 4$.
  \lname{negAllCriterion\_of\_modEq\_four} gives (Z1) or (Z2),
  \lname{hasCharPattern\_neg\_one\_of\_modEq\_four} the character~$\phi_0$, and
  \lname{hasCharPattern\_mulSingle\_one\_of\_modEq\_four} is (O2).}
\end{lemma}
\begin{proof}
  \leanok
  \uses{lem:squares2}
  We have that~$s$ is odd and $r \equiv s \bmod 4$. By assumption, (O2) holds. If~$|r|$ is not
  a square, then (Z1) holds. Otherwise, $r = \pm t^2$ with~$t$ odd, so $r \equiv \pm 1 \bmod 8$
  and $s \equiv \pm 1 \bmod 4$ (with the same sign); then $r + 2s \equiv \pm 3 \bmod 8$, so (Z2)
  holds (Lemma~\ref{lem:prime3}).
\end{proof}

One can extend these criteria to cover some further cases; see Section~\ref{sec:further}.
